\newpage
\section{angle 类}
\label{sec:class_angle}

\tkzClass{angle} 类是一个实验性的辅助对象，用于表示
由三个点定义的角。
它目前是自包含的，不与其他类交互。
其主要目的是提供一个简单直接的接口，以获取：

\begin{itemize}
  \item 有向角（弧度制），
  \item 归一化到 $[0, 2\pi)$ 的有向角，
  \item 内角（无向），
  \item 内角的度数。
\end{itemize}

\tkzClass{angle} 对象是 \emph{静态的}：所有值都在创建时计算，
之后不再更新。

\subsection{创建对象}
\label{ssub:method_angle_new}

\begin{verbatim}
  local alpha  = angle(z.A, z.B, z.C)   -- equivalent to angle:new(A,B,C)
  local beta   = angle(z.B, z.C, z.A)
  local gamma  = angle(z.C, z.A, z.B)
\end{verbatim}

三个参数分别为：
\begin{itemize}
  \item \tkzAttr{angle}{ps}：角的顶点，
  \item \tkzAttr{angle}{pa}：定义第一条射线的第一个点，
  \item \tkzAttr{angle}{pb}：定义第二条射线的第二个点。
\end{itemize}

\subsection{属性}

每个 \tkzClass{angle} 对象内部存储了以下属性：

\vspace{1em}
\bgroup
\catcode`_=12
\small
\captionof{table}{角的属性。}\label{angle:attributes}
\begin{tabular}{lll}
\toprule
\textbf{属性} & \textbf{含义} & \texttt{参考} \\
\midrule
\tkzAttr{angle}{ps}   & 角的顶点 & [\ref{par:attr_angle_ps}] \\
\tkzAttr{angle}{pa}   & 第一个定义点（射线 $[ps\,pa]$） & [\ref{par:attr_angle_pa}] \\
\tkzAttr{angle}{pb}   & 第二个定义点（射线 $[ps\,pb]$） & [\ref{par:attr_angle_pb}] \\
\tkzAttr{angle}{type} & 对象类型名，始终为 \texttt{"angle"} & \\
\tkzAttr{angle}{raw}  & 有向角（弧度），可为负 & [\ref{par:attr_angle_raw}] \\
\tkzAttr{angle}{norm} & 归一化到 $[0,2\pi)$ 的有向角 & [\ref{par:attr_angle_norm}] \\
\tkzAttr{angle}{value} & 取值范围为 $[0,\pi]$ 的无向角 & [\ref{par:attr_angle_value}] \\
\tkzAttr{angle}{deg} & 内角的度数 & [\ref{par:attr_angle_deg}] \\

\bottomrule
\end{tabular}
\egroup

\paragraph{\tkzAttr{angle}{ps}}\label{par:attr_angle_ps}
角的顶点（summit）。

\paragraph{\tkzAttr{angle}{pa}}\label{par:attr_angle_pa}
定义第一条射线 $[ps\,pa]$ 的点。

\paragraph{\tkzAttr{angle}{pb}}\label{par:attr_angle_pb}
定义第二条射线 $[ps\,pb]$ 的点。

\paragraph{\tkzAttr{angle}{raw}}\label{par:attr_angle_raw}
以弧度表示的有向角，从射线 $[ps\,pa]$ 到射线 $[ps\,pb]$，区间为 $(-\pi,\pi]$；它可以为负。

\paragraph{\tkzAttr{angle}{norm}}\label{par:attr_angle_norm}
归一化到区间 $[0,2\pi)$ 的有向角。

\paragraph{\tkzAttr{angle}{value}}
\label{par:attr_angle_value}
返回取值范围为 $[0,\pi]$ 的内角（无向）：

\texttt{示例：}
\begin{tkzexample}[latex=.45\textwidth]
  \directlua{%
  init_elements()
  z.O = point(0, 1)
  z.T = point(2, 2)
  C.OT = circle(z.O, z.T)
  z.M = C.OT:point(.13)
  A.OTM = angle(z.O, z.T, z.M)
  tkzA = A.OTM.value}
\begin{center}
\begin{tikzpicture}
  \tkzGetNodes
  \tkzDrawLines(O,T O,M)
  \tkzDrawCircle(O,T)
  \tkzDrawPoints(O,T,M)
  \tkzLabelPoints(O,T,M)
  \tkzMarkAngle(T,O,M)
  \tkzLabelAngle[pos=1.5](T,O,M){%
  \tkzPN[3]{\tkzUseLua{tkzA}}}
\end{tikzpicture}
\end{center}
\end{tkzexample}


\paragraph{\tkzMeth{angle}{deg}}
\label{par:attr_angle_deg}
返回以度数表示的内角。

\texttt{示例：}

\begin{tkzexample}[latex=.35\textwidth]
\directlua{
  init_elements()
  z.A = point(0,0)
  z.B = point(3,0)
  z.C = point(1,2)
  A.alpha = angle(z.B, z.A, z.C)
  tex.print("Angle at A = "..A.alpha.deg.." degrees")
}
\begin{center}
\begin{tikzpicture}
  \tkzGetNodes
  \tkzDrawPolygon(A,B,C)
  \tkzDrawPoints(A,B,C)
  \tkzLabelPoints(A,B)
  \tkzLabelPoints[above](C)
\end{tikzpicture}
\end{center}
\end{tkzexample}

也可以这样：

\begin{verbatim}
   T.ABC = triangle(z.A, z.B, z.C)
   local val = T.ABC.alpha_.deg
\end{verbatim}

所有值都是数值标量，对象一旦创建便保持不变。

\texttt{示例：}

\begin{tkzexample}[latex=.5\textwidth]
\directlua{
  init_elements()
  z.A = point(0, 0)
  z.C = point(3, 0)
  z.B = point(3, 2)
  T.ABC = triangle(z.A, z.B, z.C)
  A.A = T.ABC.alpha_
  A.B = T.ABC.beta_
  T.C = T.ABC.gamma_
  tex.print("A(value) = \\", A.A.value)
  tex.print('\\\\')
  tex.print("A(raw) = \\", A.A.raw)
  tex.print('\\\\')
  tex.print("A(deg) = \\", A.A.deg)}

\begin{center}
\begin{tikzpicture}
  \tkzGetNodes
  \tkzDrawPolygon(A,B,C)
  \tkzDrawPoints(A,B,C)
  \tkzLabelPoints(A,C)
  \tkzLabelPoints[above](B)
\end{tikzpicture}
\end{center}
\end{tkzexample}


\subsection{方法}

\vspace{1em}
\bgroup
\catcode`_=12
\small
\captionof{table}{angle 的方法。}\label{angle:methods}
\begin{tabular}{ll}
\toprule
\textbf{方法} & \texttt{参考} \\
\midrule
\textbf{创建} & \\
\midrule
\tkzMeth{angle}{angle(ps, pa, pb)} & [\ref{ssub:method_angle_new}] \\
\midrule
\textbf{访问器} & \\
\midrule
\tkzMeth{angle}{get()} & [\ref{ssub:method_angle_get}] \\
\midrule
\textbf{测试} & \\
\midrule
\tkzMeth{angle}{is\_direct()} & [\ref{ssub:method_angle_is_direct}] \\

\bottomrule
\end{tabular}
\egroup


\subsubsection{\tkzMeth{angle}{get()}}
\label{ssub:method_angle_get}
返回三个定义点：
\begin{verbatim}
local ps, pa, pb = alpha:get()
\end{verbatim}

\subsubsection{\tkzMeth{angle}{is\_direct()}}
\label{ssub:method_angle_is_direct}
当角为正（逆时针方向）时返回 \verb|true|。




\medskip
本类是 \emph{实验性的}，在未来版本中可能会发生变化。
